\section{Discussion}\label{sec:discussion}
The normalized smoothness coordinate yields a coherent scientific interface for the existing family of finite-difference PLS trends, while all numerical algorithms remain free to use the corresponding penalty internally. A full singular-value decomposition supplies the positive spectrum for inversion, conditional degrees of freedom, fitting, and the quantities needed by established selection scores. The shift-free expression \eqref{eq:shiftfree} changes only the arithmetic used to reconstruct the same estimator, yet its independent high-precision validation demonstrates substantially improved robustness in the tested difficult cases.

The evidence is deliberately separated into statistical and numerical questions. The RMSE pilot is an initial recovery comparison with only two true-trend shapes, two sample lengths, two noise structures, and 20 independent seeds per generating condition. It does not establish significant or universal ranking of CV, GCV, AICc and marginal ML. In particular, the observed-difference reference mean is a conditional plug-in and is not equivalent to leave-one-out refitting or joint mean estimation. The numerical comparison fixes float64 input values for $\lambda$ and $\mu$: even excellent high-precision agreement does not prove a global optimum of a nonlinear parameter search or general prediction gains.

Our boundary calculation shows that the published RSS-based BIC expression is unbounded below as the interpolating penalty is approached. Its restricted versions at $s_{\min}=0.05$ and $0.10$ are disclosed numerical policies, not theoretical replacements; developing a coherent model-based alternative is an important direction for further work.

Adopting full SVD as the framework default places a premium on verified accuracy and mathematical unity, not minimum elapsed time. Dense spectral decomposition is costly in time and memory. Weinert's specialized structured GCV methodology \citep{weinert2007} and related banded or reduced-rank methods \citep{biessy2026} should be included in a distinct future end-to-end benchmark with matched search candidates and valid outputs. An unrelated transform-based smoothing operator is not automatically a matched finite-difference baseline \citep{garcia2010}. The earlier observed benefit from cached smoothness inversion does not make $s$ intrinsically faster than $\lambda$; it amortizes work across repeated evaluations.

\section{Conclusions}
A bounded normalized smoothness parameter can be used to describe and select an established penalized least-squares trend. A full SVD of the difference operator provides an internally consistent implementation; shift-free reconstruction avoids a fragile high-degree reference and shows strong accuracy against an independent 85-digit reference in tested high-order regimes. The classical residual-based BIC expression exhibits a near-interpolation boundary divergence, requiring an explicit statistical caveat when used over a restricted grid. These results support an accuracy-oriented spectral framework and motivate larger statistical and carefully matched computational investigations.

\section*{Data and code availability}
All experiments are reproducible using the source code, immutable run-specific CSVs, numerical-validation metadata and checksum records in \url{https://github.com/heri-espino/trendestimation}, particularly the \texttt{paper/benchmarks/experiment004/} directory. The figure-generation source and Elsevier manuscript are in \texttt{paper/csda-spectral-smoothness/}. This draft does not overwrite the existing Springer manuscript.

\section*{Declaration of competing interests, funding and contributions}
These declarations and the final author ordering require coauthor confirmation prior to submission. No unverified funding or conflict statement is asserted.
